\documentclass{article}
\usepackage[UTF8]{ctex}
\usepackage{amsmath,amssymb,amsfonts}
\usepackage{enumitem}
\usepackage{geometry}
\geometry{a4paper, margin=1in}
% === 错题本排版优化（不参与题目标识） ===
\setlength{\parindent}{0pt}
\setlength{\parskip}{0.35em plus 0.10em minus 0.08em}
\setlength{\emergencystretch}{3em}
\setlist{itemsep=0.25em, topsep=0.40em, parsep=0pt}
\allowdisplaybreaks[2]
% === 错题本排版优化结束 ===
\begin{document}
\section*{哈尔滨工程大学2024年数学分析真题}
\section*{84．哈尔滨工程大学2024年数学分析试卷}

\subsection{题目1}
1．令 \(A, B\) 为任意的非空有界实数集，定义

\[
A+B=\{z \mid z=x+y, x \in A, y \in B\}, A-B=\{z \mid z=x-y, x \in A, y \in B\} .
\]

证明： \(\sup (A+B)=\sup A+\sup B, \sup (A-B)=\sup A-\inf B\) ．

\subsection{题目2}
2．计算 \(\lim _{n \rightarrow \infty} \sin (n 1 \alpha \pi)\) ，其中 \(\alpha\) 为有理数。

\subsection{题目3}
3．求曲线 \(\left(\frac{x^{2}}{a^{2}}+\frac{y^{2}}{b^{2}}\right)^{2}=x^{2}+y^{2}\) 所围平面图形的面积．

\subsection{题目4}
4．设 \(f(x, y)\) 在 \(P\left(x_{0}, y_{0}\right)\) 处可微，\(l_{1}, l_{2}, \cdots, l_{n}\) 为 \(P\left(x_{0}, y_{0}\right)\) 点处给定的 \(n\) 个方向不同的单位向量，两相邻向量夹角为 \(\frac{2 \pi}{n}\) ，计算 \(\sum_{i=1}^{n} \frac{\partial f\left(x_{0}, y_{0}\right)}{\partial l_{i}}\) ，其中 \(\frac{\partial f\left(x_{0}, y_{0}\right)}{\partial l_{i}}\) 表示 \(f(x, y)\) 在 \(P\left(x_{0}, y_{0}\right)\) 处沿 \(l_{i}\) 的方向导数。

\subsection{题目5}
5．令 \(x=\frac{1}{n} \sum_{k=1}^{n} x_{k}\) ，其中 \(x_{i} \in(0, \pi)\) ，证明：\(\prod_{k=1}^{n} \frac{\sin x_{k}}{x_{k}} \leq\left(\frac{\sin x}{x}\right)^{n}\) ．

\subsection{题目6}
6 ．求 \(f(x)=[x] \sin \pi x\) 的单侧导数．

\begin{itemize}

\subsection{题目7}
\end{itemize}

\subsection{题目7}
原题边界“7．”：如果 \(f(x)\) 在 \([0,1]\) 上二次可微，\(f^{\prime \prime}(x)\) 有界且黎曼可积．证明：
\[
\lim _{n \rightarrow \infty} n^{2}\left(\int_{0}^{1} f(x) \mathrm{d} x-\frac{1}{n} \sum_{i=1}^{n} f\left(\frac{2 i-1}{2 n}\right)\right)=\frac{f^{\prime}(1)-f^{\prime}(0)}{24} .
\]

\subsection{题目8}

\subsection{题目8}
原题边界“8．”：令
\[
u_{n}(x)= \begin{cases}-n, & x \leq-n \\ x, & -n<x \leq n \\ n, & x>n\end{cases}
\]
\(f(x)\) 为实值函数，证明：\(f(x)\) 连续等价于对任意固定的正整数 \(n, g_{n}(x)=u_{n}(f(x))\) 都是 \(x\) 的连续函数．

\subsection{题目9}
9．假设 \(f(x) \in C^{\infty}([0,1])\) 且满足下列条件：

\begin{itemize}
  \item[（1）] \(f(x)\) □ 0.
  \item[（2）] \(f^{(n)}(x)=0, n=0,1,2, \cdots\)
  \item[（3）] 实数序列 \(\left\{a_{n}\right\}\) 满足级数 \(\sum_{n=1}^{\infty} a_{n} f^{(n)}(x)\) 在 \([0,1]\) 上一致收敛。
\end{itemize}

证明： \(\lim _{n \rightarrow \infty} n!a_{n}=0\) ．

\subsection{题目10}
10．证明：\(H_{n}(x)=(-1)^{n} e^{x^{2}} \frac{\mathrm{~d}^{n}}{\mathrm{~d} x^{n}}\left(e^{-x^{2}}\right)\) 所有根都是实数．

\subsection{题目11}
11．计算第二类曲面积分

\[
\iint_{S}(f(x, y, z)+x) \mathrm{d} y \mathrm{~d} z+2 f(x, y, z) \mathrm{d} z \mathrm{~d} x+(f(x, y, z)+z) \mathrm{d} x \mathrm{~d} y .
\]

其中 \(f(x, y, z)\) 为连续函数，\(S\) 为平面 \(x-y+z=1\) 在第四卦限的上侧．

\subsection{题目12}
12．设 \(f(x, y)\) 定义在 \(x^{2}+y^{2} \leq 1\) ，且具有连续的偏导数，\(|f(x, y)| \leq 1\) 。证明：在 \(x^{2}+y^{2} \leq 1\) 内有一点 \(\left(x_{0}, y_{0}\right)\) ，使得 \(\left|f_{x}\left(x_{0}, y_{0}\right)\right|^{2}+\left|f_{y}\left(x_{0}, y_{0}\right)\right|^{2}<16\) ．
\end{document}
